\documentclass[10pt,a4paper]{article}
\usepackage[margin=22mm]{geometry}
\usepackage{amsmath,amssymb,booktabs,longtable,array}
\usepackage[T1]{fontenc}
\usepackage{hyperref}
\newcommand{\Lean}[1]{\texttt{\detokenize{#1}}}
\newcommand{\bytecode}{\textsf{bytecode fact}}
\newcommand{\derived}{\textsf{derived from bytecode}}
\newcommand{\model}{\textsf{certified model assumption}}
\newcommand{\runtime}{\textsf{runtime contract}}
\newcommand{\obligation}{\textsf{future Lua implementation obligation}}
\title{Geometry- and Physics-Guided IR Train Proof Contract}
\author{Immersive Railroading 1.11.0 / Minecraft 1.7.10}
\date{2026-08-16}

\begin{document}
\sloppy
\setlength{\tabcolsep}{2pt}
\setlength{\emergencystretch}{3em}
\hfuzz=100pt
\hbadness=10001
\maketitle

\begin{center}
\fbox{\parbox{0.90\linewidth}{\centering
\textbf{NOT\_READY}. This is a proof and implementation-contract artifact.
The production Lua controller is unchanged. All previous PID proof,
optimality, and readiness claims are withdrawn.}}
\end{center}

\tableofcontents

\section{Authority, scope, and classifications}

The authoritative specifications are \texttt{PLAN\_NEWPROOF.md} and its
overriding supplement \texttt{PLAN\_NEWPROOF\_REVISION.md}. The proof target
is a robust stopping and terminal-arrival contract for a directed station or
signal. It is not a PID proof and it does not claim controller optimality.

The source-connected funnel is:
\[
 \begin{gathered}
 \text{tracer samples}\to\text{track-centre splines}\to\text{routing graph}\to\\
 \text{detector profile}\to\text{front coupler}\to\text{combined command}\to\\
 \text{IR plant}\to\text{arrival/holding}.
 \end{gathered}
\]
The spline is geometric source of truth. The graph is only a routing
abstraction and never fits or rewrites a spline. Every statement is labelled
as one of: proved by Lean; \bytecode; \derived; \model; \runtime;
\obligation; or unsupported/withdrawn.

\section{Provenance and Java-8 source boundary}

The extracted directory \texttt{immersive\_railroading/.cache/jar/} is
authoritative. Its manifest identifies Minecraft 1.7.10, Immersive Railroading
1.11.0, and \texttt{Multi-Release: true}. Root class files have major version
52 (Java 8); \texttt{META-INF/versions/16} classes have major version 60 and
are excluded. Root class and complete decompiler-input hashes are recorded in
\texttt{pid-jar-manifest.md}; the six-input reproducibility check is recorded
in \texttt{pid-physics-contract-manifest.sha256}. The exact branch ledger is
in \texttt{pid-bytecode-traceability.md}.

The root \texttt{javap -p -c} evidence establishes the following boundaries:
\begin{itemize}
 \item \texttt{CommonAPI.info()} provides dynamic identity, definition,
 weight, speed, direction, brake/control, and locomotive fields;
 \texttt{consist()} is aggregate validation only.
 \item The detector manager emits \texttt{ir\_train\_overhead} with a stock
 UUID, and detector \texttt{info()} delegates through the API.
 \item \texttt{EntityRollingStock} position is not proved to be the F3
 hitbox or rail centreline. \texttt{SimulationState.Configuration} and
 \texttt{calculateCouplerPositions()} consume bogey, coupler, yaw, track,
 mass, slack, and resistance data.
 \item \texttt{CubicCurve.position()} is a cubic Bernstein expression;
 \texttt{lengthWithCache()} samples derivatives and accumulates trapezoids.
 \item \texttt{SimulationState.forcesNewtons()} computes signed grade plus
 tractive effort. \texttt{frictionNewtons()} computes the resistance and
 brake branches listed in Section~\ref{sec:physics}.
 \item \texttt{SimulationState.next()} clones state, calls movement, has an
 unchanged-position velocity-zero branch, recomputes couplers/collisions and
 direct resistance. \texttt{Consist\$Particle} updates force, friction,
 linkage, position, and velocity per particle.
\end{itemize}

\section{Trace-centre and spline model}

\subsection{All reconstruction variables}

The sample index is $i\in\mathbb{N}$. Each accepted tracer observation is
\[
 o_i=(t_i,r_i,u_i,h_i,\mathrm{uuid}_i,\mathrm{def}_i,
      b_i,\alpha_i,\Delta\ell_i,\kappa_i,\epsilon_{float,i},
      \epsilon_{missing,i},e_i),
\]
where $t_i$ is the tick, $r_i\in\mathbb{R}^3$ is raw \texttt{getPos()},
$u_i$ is speed, $h_i$ is heading/orientation, $b_i$ is the compact tracer
bogey/reference offset, $\alpha_i\geq0$ is observation age,
$\Delta\ell_i\geq0$ is sampling spacing, $\kappa_i\geq0$ is a curvature
bound, $\epsilon_{float,i}\geq0$ is numeric error,
$\epsilon_{missing,i}\geq0$ is missing/delayed-observation error, and
$e_i\in\mathbb{R}^3$ is the measured residual. The centre estimate and bound
are
\[
 c_i=r_i+b_i,\qquad
 \epsilon_{trace,i}=\Delta\ell_i\kappa_i+epsilon_{float,i}
                         +\epsilon_{missing,i},qquad
 c_i^{actual}=c_i+e_i,\quad\|e_i\|_1\leq\epsilon_{trace,i}.
\]
The front coupler is absent from this record. It is a stopping reference only.
The offset, validity, and residual are a \model{} and \runtime{} boundary
until the future sampler supplies them; raw position is never silently reused.

For stable spline ID $j$, define four controls
$P_{j,0},P_{j,1},P_{j,2},P_{j,3}\in\mathbb{R}^3$ and
\[
 \sigma_j(t)=(1-t)^3P_{j,0}+3(1-t)^2tP_{j,1}
       +3(1-t)t^2P_{j,2}+t^3P_{j,3},\quad 0\leq t\leq1.
\]
The stored variables are
\[
 S_j=(j,P_{j,0:3},A_j,L_j,\tau_j,g_j,\kappa_j,
      \epsilon_{fit,j},T_j,M_j),
\]
where $A_j$ is a monotone arc-length table, $L_j>0$ is total length,
$\tau_j$ is tangent, $g_j$ is grade, $\kappa_j$ is curvature or a bound,
$\epsilon_{fit,j}\geq0$ is fitting error, $T_j$ is the list of source trace
IDs, and $M_j$ is optional builder metadata. Straight track is a degenerate
cubic. IR turns are not called exact circles: the root turn implementation
uses a cubic approximation.

For each observation, the fit stores a parameter $q_i\in[0,1]$ and requires
\[
 \|c_i-\sigma_j(q_i)\|_1
 \leq\epsilon_{fit,j}+\epsilon_{trace,i}.
\]
The producer \Lean{validatedSplineProducer} returns a validated spline only
when IDs, length, fit error, source traces, all observations, parameter rows,
and residual rows pass. Thus the theorem does not accept an arbitrary caller
certificate as proof of its own validity.

\subsection{Incremental reconstruction and topology}

For a new trace and existing spline, the match record is
\[
 m=(d_3,d_\tau,d_\kappa,d_s;\delta_3,\delta_\tau,
    \delta_\kappa,\delta_s,h,\delta_h),
\]
where the first four values are 3-D position, tangent, curvature, and
arc-coordinate discrepancies, and the deltas are certified thresholds. A
match requires all four 3-D/arc inequalities; 2-D projection intersection is
never enough. The topology class is one of continuation, extension, split,
overlap, shared connection, isolated crossing, overpass, new geometry, or
unresolved. A same-level pair whose paths both continue without connection is
an isolated crossing. A height separation greater than tolerance is an
overpass. Both preserve separate paths and are not routable. Unresolved
branches remain unavailable; manual classification is an explicit later
transition, never a silent guess.

The implementable state transition is:
\begin{enumerate}
 \item validate observations and construct $c_i$ and its residual terms;
 \item fit and validate the cubic, arc marks, reverse controls, length,
 tangent, grade, curvature, and fit error;
 \item compare against existing splines using all 3-D/arc metrics;
 \item classify the topology, preserving unresolved/crossing/overpass state;
 \item commit only validated spline/source metadata and revision;
 \item rebuild the graph from validated spline IDs and explicit connections.
\end{enumerate}
This is the \model{} and \obligation{} for the future persistent mapper.

\section{Routing graph, stations, signals, and active routes}

At graph revision $r_G$, define $G_{r_G}=(V_{r_G},E_{r_G})$. A vertex is
$v=(id,j,s)$ with $0\leq s\leq L_j$. A directed edge is
\[
 e=(id_e,j,s_0,s_1,o,v_a,v_b),\qquad o\in\{+1,-1\},
\]
referencing a validated spline ID and arc interval. A physical spline normally
produces symmetric forward and reverse edges. A connection record is explicit
evidence between validated endpoints; it is not inferred from a projected
crossing.

A marker is
\[
 m=(id_m,p_m,\delta_m,a_m,k_m),
\]
where $p_m$ is its world point, $\delta_m$ is mapping tolerance, $a_m$ is
exactly one permitted approach orientation, and $k_m$ is station or signal
kind. The closest spline candidate is accepted only when distance is within
$\delta_m$ and $0\leq s_m\leq L_j$. The opposite approach requires a separate
marker configuration.

An active route is $R=(r_G,S_R,V_R,E_R,O_R)$. It validates only its referenced
splines, vertices, edges, intervals, endpoint sequence, orientation, and the
current graph revision. It does not scan the whole graph. The graph-network
producer consumes candidate traces, invokes the validated spline producer for
each candidate, and rejects empty or duplicate stable identifiers rather than
silently aliasing geometry. A changed active section
invalidates $R$, invokes route finding against the latest graph, and fails
closed with a report if no route exists. The Lean producers are
\Lean{directedMarkerProducer}, \Lean{routingGraphProducer},
\Lean{graphNetworkProducer}, \Lean{routeReferencesGraph}, and
\Lean{routeDecision}.

\section{Detector profile and source-separated data}

For each \texttt{ir\_train\_overhead} event, record $u_k=$\texttt{stock\_uuid}
and immediately call detector \texttt{info()}. The ordered record is
\[
 q_k=(u_k,d_k,w_k,v_k,a_k,b^{train}_k,b^{ind}_k,T_k,C_k,
      \eta^{sys}_k,\eta^{adh}_k,\mu^{shoe}_k),
\]
containing UUID, definition ID, weight, direction, speed, both brake channels,
traction $T$, cogging $C$, brake-system and adhesion efficiencies, and shoe
friction when available. Dynamic \texttt{info()} data, static definition data,
global configuration, and runtime certificates are separate records.
\texttt{consist()} validates aggregate information only. The catalogue is
immutable during a run; UUID/order/definition/weight/brake mismatch reports,
stops, and requires recataloguing. Initial catalogue speed is approximately
10 km/h; a higher value requires a bound derived from detector length,
minimum spacing, event latency, and \texttt{info()} response time.

The Lean boundary is \Lean{catalogueProducer}\,$\to$\Lean{StockRecord}
$\to$\Lean{FrozenProfile}\,$\to$\Lean{sourcePhysicsRecordFromProfile}.
Missing info, missing static fields, missing global fields, and definition
mismatch return \texttt{none}; no neutral physical default is introduced.

\section{Directed front-coupler coordinate}

Choose the spline orientation permitted by the station/signal marker and let
$s$ increase in that direction. If $s_{fc}$ is the front-coupler coordinate,
then
\[
 x=s_{target}-s_{fc},\qquad v_s=\frac{d s_{fc}}{dt}.
\]
The stopping theorem is entered only for $x>0$ and $v_s>0$. Wrong-way motion,
target crossing, stale data, invalid localization, and unavailable route
orientation are recovery or fail-closed states.

If $s_c$ is the tracer-centre coordinate and $\delta_{fc}$ is the signed
definition/coupler offset in the selected orientation, the localization record
is
\[
 (s_c,\delta_{fc},s_{fc}=s_c+\delta_{fc},\epsilon_{fc}).
\]
The producer \Lean{frontCouplerLocalizationProducer} consumes a validated
spline, a complete coupler definition, and a centre error. Raw \texttt{getPos()}
and a model hitbox are not substitutes. Reverse orientation reverses the
directed velocity and selects reverse offsets.

\section{IR physics derivation}\label{sec:physics}

\subsection{Source records and constants}

The exact source record contains current mass $m$, design mass $m_d$, pitch
$\theta$ and separately supplied $\sin(\theta)$, slope multiplier $M_s$,
tractive effort $T$, rolling coefficient $c_r$, speed-retarder resistance
$R_s$, direct coefficient $c_d$, interference $I$, block hardness $H$, train
pressure $p_t$, independent pressure $p_i$, brake-system efficiency $\eta_s$,
adhesion efficiency $\eta_a$, global multiplier $\beta$, velocity $v$, shoe
coefficient $\mu_{shoe}$, and cogging flag. The Java bytecode evaluates the
sine; the rational model requires the runtime adapter to provide it rather
than replacing it by zero.

The source constants are
\[
 \mu_{steel}^{static}=0.70=\frac7{10},\quad
 \mu_{steel}^{kinetic}=0.42=\frac{21}{50},\quad
 \mu_{steel/castiron}^{kinetic}=0.25=\frac14,
 \quad g=9.8=\frac{49}{5}.
\]
The cast-iron value is recorded as a material-definition fact and source
field. The inspected \texttt{SimulationState.frictionNewtons()} wheel-slip
branch invokes steel/steel kinetic friction, so the proof does not silently
substitute $0.25$ into that branch. Cogging overrides for brake-system and
adhesion efficiency are recorded as bytecode facts, not inferred defaults.

\subsection{Line-by-line force ledger}

The root bytecode first gives signed grade and traction:
\[
 F_g=m(-g)\sin(\theta)M_s,
 \qquad F_T=T.
\]
The configuration constructor derives
\[
 A_d=m_d\,\mu_{steel}^{static}\,g\,\eta_s,
 \qquad A_{max}=m\,\mu_{steel}^{static}\,g\,\eta_a.
\]
The friction method then performs, in order,
\begin{align*}
 p&=\min(1,\max(p_t,p_i)),\\
 F_{candidate}&=A_d p,\\
 F_{roll}&=c_r m g,\\
 F_{zero}&=\begin{cases}0.001mg&v=0,\\0&v\ne0,\end{cases}\\
 F_{int}&=I\,1000\,H,\\
 F_{direct}&=R_s+c_d p_i m g+c_d p_t m g,\\
 F_{shoe}&=\begin{cases}m\mu_{steel}^{kinetic}&F_{candidate}>A_{max}\land|v|>0.01,\\F_{candidate}&\text{otherwise},\end{cases}\\
 F_{brake}&=\beta F_{shoe},\\
 F_{net}&=F_g+F_T-(F_{roll}+F_{zero}+F_{direct}+F_{int}+F_{brake}).
\end{align*}
The root wheel-slip branch also sets its sliding state. API command values are
validated and normalized at a separate boundary; the root physical pressure
branch itself is the upper clamp shown above. This distinction preserves both
the source equation and fail-closed command conversion.

The Lean structure \Lean{ExactForceLedger} stores every summand and source
reference. \Lean{source\_force\_ledger\_equation} proves the signed sum and
\Lean{source\_wheel\_slip\_changes\_branch} proves the slip branch. A
missing source field returns \texttt{none}; no field is replaced with zero.

\subsection{Particle aggregation}

For each particle $k$, the source-shaped transition stores state, interacting
mass $m_k$, interacting friction $f_k$, previous/next linkage, push/pull
capability, slack, collision observation, and track-transition observation.
The selected friction is direction-aware and capped by force magnitude; the
model update is
\[
 a_k=\frac{F_k-\operatorname{signedFriction}_k}{m_k},\quad
 s'_k=s_k+v_k\Delta t,\quad v'_k=v_k+a_k\Delta t,
\]
followed by linkage distance correction. \Lean{sourceParticleTransition}
rejects invalid mass/linkage and \Lean{sourceConsistTransition} preserves
per-particle order. This is a conservative source-shaped model; complete
Java collision, track lookup, pressure propagation, and all linkage branches
remain a refinement obligation.

\section{Combined command and plant transition}

The command is
\[
 u=(u_T,u_{train},u_{ind},u_E,\alpha_{cmd},\alpha_{pressure}),
\]
with throttle, train brake, independent brake, emergency state, command age,
and pressure age. The controller path is explicitly
\[
 \begin{aligned}
 &\text{spline localization}\to(x,v_s)\to
 (m,L,\kappa,\text{slack},\text{push/pull})\to u_T\\
 &\to(u_{train},u_{ind},u_E)\to F_{net}\to\text{particle step}
 \to\text{next state}.
 \end{aligned}
\]
Positive throttle is not assumed zero after stopping commitment. It is included
in the positive-traction upper bound. \Lean{apiCommandProducer} rejects
nonfinite and out-of-range raw values, \Lean{exactCommandPhysics} constructs
the command-modified exact ledger, and \Lean{exactPlantTransition} applies
that ledger to a discrete state transition. \Lean{api\_command\_to\_exact\_plant}
connects the raw API producer to existence of a next state. Invalid command or
linkage input is \texttt{none}; this is the fail-closed path.

Saturation and slew are represented as explicit command/runtime fields to be
implemented and measured. They may not be omitted because they are inconvenient.

\section{Uniform guard, safety, arrival, and holding}

Let $B_{min}$ be a lower brake-force bound valid over every declared source
branch and profiled particle. It must already account for rolling/direct/
interference/near-zero behavior, wheel slip, material/efficiency branches, and
consist aggregation. Let $T_{max}$ bound positive traction, and let
$G_{max},K_{max},S_{max},P_{max}$ bound adverse grade, curvature, slack, and
push/pull forces. For mass $m>0$,
\[
 a_{min}=\frac{B_{min}-T_{max}-G_{max}-K_{max}-S_{max}-P_{max}}{m}>0.
\]
The delay budget is
\[
 \tau=\tau_{sample}+\tau_{command}+\tau_{pressure}
       +\tau_{detector}+\tau_{route},
\]
where jar-internal next-branch timing is deterministic source evidence and
the other terms are measured runtime contracts. For distance $d$, speed $v$,
localization error $e_g$, and terminal tolerance $e_t$, the terminal-entry
guard is
\[
 d>0,\quad v>0,\quad a_{min}>0,\quad
 d\geq v\tau+\frac{G_{max}+K_{max}+S_{max}+P_{max}}{2m}\tau^2
       +\frac{v^2}{2a_{min}}+e_g+e_t.
\]
It includes delay distance, adverse delayed motion, stopping distance,
geometry error, and terminal tolerance. \Lean{terminalGuard} rejects wrong
way, crossed target, and nonpositive-deceleration inputs.

The safety envelope is $d\geq0$, $v\geq0$, and $v^2\leq2a_{min}d$.
\Lean{arrival\_step\_preserves\_no\_crossing} and
\Lean{arrival\_run\_velocity\_nonincreasing} prove the explicit rational
discrete claims. Terminal arrival is separate:
\Lean{terminal\_arrival\_step} and
\Lean{terminal\_arrival\_step\_position} prove $v_s=0$ and $x=0$ for the
certified terminal step. Runtime acceptance additionally requires
\[
 |v_s|\leq\epsilon_v,\qquad |x|\leq\epsilon_{stop},
 \qquad\text{stable for }N_{stable}\text{ observations}.
\]
The position tolerance is derived, not chosen for convenience:
\[
 \epsilon_{stop}\geq\epsilon_{trace}+\epsilon_{fit}+epsilon_{coupler}
 +\epsilon_{sample-age}+\epsilon_{command-delay}+\epsilon_{numeric}.
\]
Holding uses the actual IR condition
$v_s=0\land|F|<F_{friction}$, not merely zero command or zero force.
\Lean{holdingTransition} preserves position and zero velocity under that
condition. Stable acceptance is a runtime contract, not a certificate that
the unchanged Lua already observes the required window.

\section{Producer/consumer theorem inventory}

\tiny
\begin{longtable}{>{\raggedright\arraybackslash}p{0.23\linewidth}
                  >{\raggedright\arraybackslash}p{0.30\linewidth}
                  >{\raggedright\arraybackslash}p{0.40\linewidth}}
\toprule
Arrow & Producer and consumer & Lean relation and classification \\
\midrule
Trace $\to$ spline & observations/residuals $\to$ validated spline &
\Lean{validated\_spline\_trace\_fit\_is\_checked}; proved by Lean under trace model \\
Spline $\to$ graph & validated spline/candidates $\to$ directed edge/network &
\Lean{validated\_spline\_to\_directed\_routing\_edge},
\Lean{graph\_network\_consumes\_only\_source\_validated\_splines}; proved by Lean \\
Graph/profile $\to$ localization & route/marker/definition $\to$ front coupler &
\Lean{marker\_maps\_to\_valid\_spline\_interval}, \Lean{coupler\_localization\_consumes\_validated\_spline}; proved by Lean \\
Profile $\to$ physics & dynamic/static/config/runtime $\to$ exact source input &
\Lean{source\_profile\_physics\_preserves\_mass\_and\_channels}; derived/ proved \\
Command $\to$ plant & numeric/API command $\to$ exact ledger/next state &
\Lean{api\_command\_to\_exact\_plant}, \Lean{exact\_plant\_has\_source\_next\_state}; conditional model \\
Plant $\to$ arrival & force/transition $\to$ guard, safety, terminal acceptance &
\Lean{arrival\_step\_preserves\_no\_crossing}, terminal theorems, stable contract; conditional model \\
\bottomrule
\end{longtable}
\normalsize

\subsection{Required spline-to-edge proof}

The theorem \Lean{validated\_spline\_to\_directed\_routing\_edge} takes the
producer equality
\Lean{validatedSplineProducer c = some v}, orientation $o$, parameter $t$,
error $\epsilon$, point $q$, interval premise $0\leq t\leq1$, and pointwise
premise $\|q-\sigma_c(t)\|_1\leq\epsilon$. It proves the same stable spline
ID, both forward/reverse endpoint correspondences, the directed coordinate
mapping, point correspondence through the reverse cubic, source trace IDs,
and preservation of the localization inequality. It depends on
\Lean{traceFitRowsValid}, \Lean{splineCandidateValid},
\Lean{validatedSplineProducer}, \Lean{cubicReverse},
\Lean{cubic\_reverse\_position}, \Lean{edgeCoordinate},
\Lean{directedCoordinateOfParameter}, \Lean{edgePointAt}, and
\Lean{directedEdge}. Its transitive \texttt{\#print axioms} output is
\texttt{[propext, Classical.choice, Quot.sound]}. The tests mutate direction,
interval, fit row, and spline ID.

\section{Anti-cheating audit and gate result}

The audit applies the thirteen base anti-cheating classes and the revision's
additional checks: no forbidden proof escape; no vacuous producer; no
specialized example promoted to a theorem; no equation mismatch; no
disconnected definition; no controller-shaped optimality; no sign/dimensional
error; no temporal under-approximation; no API/refinement alias; no
provenance drift; no rational/floating conflation; no safety/liveness scope
confusion; and meaningful negative/mutation tests. Additional audit checks
reject optimistic defaults, self-authenticated certificates, and fake
transition aliases.

Gates A--H pass only at the explicit proof/contract boundary and are marked
conditional where live source or timing evidence is required. Gate I is
\textbf{NOT\_READY}: full Lua source adapters, measured external timing,
complete Java linkage/collision refinement, and a live mutation harness remain
future obligations. No gate is marked passed merely because Lean or LaTeX
compiles.

\section{Implementation contract and readiness}

The future Lua implementation must persist trace/spline/topology records,
derive the graph only from validated splines, map directed markers, catalogue
each detector stock immediately, freeze and validate the profile, localize the
front coupler, send all command channels together, certify the complete
uniform force/delay bound, implement stable terminal holding, and fail closed
with diagnostics on every missing or mutated input. It must run the vectors
against curved, slack, push/pull, adhesion-limited, wheel-slip, detector-delay,
route-mutation, and profile-mutation scenarios.

The final classification is:
\begin{itemize}
 \item Lean relations and exact vectors: \textbf{proved by Lean} under their
 explicit premises;
 \item root equations, constants, versions, classes, and branch locations:
 \textbf{derived from bytecode} or bytecode facts;
 \item trace residual, arc-length approximation, physical force bounds,
 curvature/slack/push-pull bounds: \textbf{certified model assumptions};
 \item detector delivery, profile immutability, numeric conversion, timing,
 route revision, diagnostics, and fail-closed behavior: \textbf{runtime
 contracts};
 \item persistent geometry/profile/controller/plant integration and live
 mutation tests: \textbf{future Lua implementation obligations};
 \item PID optimality, exact floating-point correctness, complete Java
 linkage/collision identity, and deployment approval: \textbf{unsupported or
 withdrawn}.
\end{itemize}

\begin{center}
\fbox{\textbf{NOT\_READY -- not approved for user deployment}}
\end{center}

\end{document}
